\subsection{KLM preferential representation: the general construction}
\label{sec:p08-klm-preferential-representation}

The preferential theorem is a representation of a whole relation between
formulas, not an assertion that a particular finite default example is
adequate. We reconstruct the canonical copied-state construction of
Kraus, Lehmann, and Magidor. The page locators below refer to their freely
offered author posting, \texttt{cs/0202021v1}, pp.~6--7, 10--19 and
25--34; its 1990 journal identity remains distinct
\cite{KrausLehmannMagidor1990,KrausLehmannMagidor2002Author}.
The earlier statement in Theorem~\ref{thm:p08-klm-representation} uses
the same preferential relation, but its cited standing is not the proof
of the construction supplied here.

\begin{hbdefinition}{def:p08-klm-language}{Fixed classical language and universe of reference}
Let \(L\) be a set of finitary classical propositional formulas, closed
under the classical connectives, with \(\neg,\vee\) primitive and
\(\wedge,\to,\leftrightarrow\) defined classically. Let \(U\) be a
set of classical valuations, possibly restricted by a background theory.
Write \(m\models\alpha\) for satisfaction, \(\models\alpha\) for
truth at every \(m\in U\), and \(\alpha\models\beta\) for
\(\models(\alpha\to\beta)\). These are the source's relative
notations: \(\to\) is material implication and \(\models\) refers
to the fixed \(U\), not to the new preferential relation.
Assume semantic compactness for \((L,U)\): every set of formulas whose
finite subsets have models in \(U\) has a model in \(U\).
This is an explicit source hypothesis, not a conclusion of preferential
representation. Ordinary propositional valuations, including the models
of a fixed background theory, supply the usual compact setting.
\end{hbdefinition}

\begin{hbdefinition}{def:p08-klm-assertions}{Flat conditional assertions and System P}
An assertion \(\alpha\sim\beta\) is a pair in \(L\times L\),
with antecedent \(\alpha\) and consequent \(\beta\). It is not a
formula constructor that can be nested under \(\neg\) or \(\vee\).
A relation \(\sim\subseteq L\times L\) is cumulative (System C)
when closed under Reflexivity, Left Logical Equivalence, Right Weakening,
Cut and Cautious Monotonicity. It is preferential (System P) when also
closed under Or. This is the native basis of Definitions~1 and~15 of
the source, not an added monotonic-premise rule. The host notation
\(\alpha\mathrel{|\!\sim}\beta\) denotes this same kind of relation.
For a set \(K\subseteq L\times L\), possibly infinite,
\(K\vdash_P\alpha\sim\beta\) means a finite derivation using the
following sheet, with assertions of \(K\) as available initial lines.
No uniform-substitution rule is included: fixed background-validity
conditions need not be invariant under substitution.
\end{hbdefinition}

\begin{hbdefinition}{def:p08-klm-model}{Smooth preferential models, with state copies}
A model is \(W=\langle S,l,\prec\rangle\), where \(S\) is a set
of states, \(l:S\to U\), and \(\prec\) is a strict partial order
(irreflexive and transitive), oriented from preferred to less preferred.
Distinct states may have the same label, and \(l\) need not cover all
of \(U\). Put
\[
 \widehat\alpha=\{s\in S:l(s)\models\alpha\},\qquad
 \mu(X)=\{s\in X:\nexists t\in X\ (t\prec s)\}.
\]
The model is smooth when, for every \(\alpha\in L\) and
\(s\in\widehat\alpha\), either \(s\in\mu(\widehat\alpha)\)
or some \(t\in\mu(\widehat\alpha)\) satisfies \(t\prec s\).
Its relation is \(\alpha\sim_W\beta\) iff
\(\mu(\widehat\alpha)\subseteq\widehat\beta\).
Empty antecedent sets give vacuous consequence; the empty state set is
allowed. Smoothness is required on all formula-definable sets, not merely
on the premises in one knowledge base, and not on every arbitrary subset
unless that stronger condition is separately imposed.
\end{hbdefinition}

% KLM native basis: C (five rules) plus Or. The final three rules are derived.
\calculussheet{sheet:p08-klm-preferential}{KLM Preferential System P}{Language: flat assertions $\alpha\sim\beta$ with $\alpha,\beta\in L$; $\to$ is classical material implication, not $\sim$; classical validity $\models$ is relative to the fixed compact universe $U$; derivations from $K\subseteq L\times L$ are finite sequences; negation is classical inside $L$; equality is not an added object-language connective; no uniform-substitution rule under fixed background constraints.}
\hbrule{rule:p08-klm-ref}{Reflexivity}{}{\alpha\sim\alpha}{\none}
\hbrule{rule:p08-klm-lle}{Left Logical Equivalence}{\models(\alpha\leftrightarrow\beta)\qquad\alpha\sim\gamma}{\beta\sim\gamma}{validity is relative to the same fixed universe $U$}
\hbrule{rule:p08-klm-rw}{Right Weakening}{\models(\beta\to\gamma)\qquad\alpha\sim\beta}{\alpha\sim\gamma}{validity is relative to the same fixed universe $U$}
\hbrule{rule:p08-klm-cut}{Cut}{\alpha\sim\beta\qquad\alpha\wedge\beta\sim\gamma}{\alpha\sim\gamma}{\none}
\hbrule{rule:p08-klm-cm}{Cautious Monotonicity}{\alpha\sim\beta\qquad\alpha\sim\gamma}{\alpha\wedge\beta\sim\gamma}{\none}
\hbrule{rule:p08-klm-or}{Or}{\alpha\sim\gamma\qquad\beta\sim\gamma}{\alpha\vee\beta\sim\gamma}{\none}
\hbrule{rule:p08-klm-sc}{Supraclassicality}{\alpha\models\beta}{\alpha\sim\beta}{derived from Reflexivity and Right Weakening; see Annex~\ref{annex:p08.klm.1}}
\hbrule{rule:p08-klm-and}{And}{\alpha\sim\beta\qquad\alpha\sim\gamma}{\alpha\sim\beta\wedge\gamma}{derived from the native cumulative rules; see Annex~\ref{annex:p08.klm.1}}
\hbrule{rule:p08-klm-s}{S}{\alpha\wedge\beta\sim\gamma}{\alpha\sim\beta\to\gamma}{derived using Or; it does not authorize the converse strengthening rule}


\begin{hbtheorem}{thm:p08-klm-basis}{proved=annex:p08.klm.1}{The native and expanded bases, with needed derived rules}
The native six-rule P basis derives Supraclassicality and And. Conversely
the seven preferential principles in Table~\ref{tab:p08-system-p-rules}
derive the native basis, without Rational Monotony.
The following are also derived: antecedents mutually consequential under
\(\sim\) have the same consequences; \(\alpha\wedge\beta\sim\gamma\)
implies \(\alpha\sim\beta\to\gamma\); and
\(\alpha\sim\gamma,\beta\sim\delta\) imply
\(\alpha\vee\beta\sim\gamma\vee\delta\).
With And and Or available, the displayed implication rule derives Cut.
All these derivations use validity relative to the same \(U\).
\end{hbtheorem}

\begin{hbtheorem}{thm:p08-klm-soundness}{proved=annex:p08.klm.2}{Soundness of every preferential rule}
Every smooth preferential model validates Reflexivity, Left Logical
Equivalence, Right Weakening, Cut, Cautious Monotonicity, Or,
Supraclassicality and And. In fact, if \(\alpha\sim_W\beta\),
\[
 \mu(\widehat\alpha)=\mu(\widehat{\alpha\wedge\beta}).
\]
Thus adding an expected conclusion to the antecedent preserves every
consequence: this is cumulativity, not arbitrary left strengthening.
Smoothness is essential to the reverse minimal-set inclusion and to
Cautious Monotonicity; the proof of Cut does not itself need smoothness.
\end{hbtheorem}

\begin{hbdefinition}{def:p08-klm-normal}{Normal worlds for a formula}
For a preferential relation define
\(C(\alpha)=\{\gamma\in L:\alpha\sim\gamma\}\) and
\(N(\alpha)=\{m\in U:\forall\gamma\in C(\alpha),\ m\models\gamma\}\).
A world in \(N(\alpha)\) is \(\alpha\)-normal. Reflexivity makes
it satisfy \(\alpha\). No existence of such a world is included in
this definition; that requires the next compactness argument.
\end{hbdefinition}

\begin{hblemma}{thm:p08-klm-normal-witness}{proved=annex:p08.klm.3}{The normal-world witness, including infinite theories}
\(\alpha\sim\beta\) iff every \(\alpha\)-normal world satisfies
\(\beta\). If \(\alpha\not\sim\beta\), there is an
\(m\in N(\alpha)\) with \(m\not\models\beta\).
The proof uses finite And closure and Right Weakening to establish
finite satisfiability of \(C(\alpha)\cup\{\neg\beta\}\), and then
the stated compactness of \(U\); it is not restricted to finite \(C(\alpha)\).
\end{hblemma}

\begin{hbdefinition}{def:p08-klm-ordinary}{The source's formula preorder}
Write \(\alpha\leq\beta\) iff \(\alpha\vee\beta\sim\alpha\).
This is the source's ``not less ordinary'' relation, not material
entailment or the preference order between states. In particular
\(\alpha\vee\beta\leq\alpha\) and
\(\alpha\vee\beta\leq\beta\).
\end{hbdefinition}

\begin{hblemma}{thm:p08-klm-preorder}{proved=annex:p08.klm.4}{Preorder and normal-world transport}
The relation \(\leq\) is reflexive and transitive.
If \(\alpha\leq\beta\), \(\beta\sim\delta\), then
\(\alpha\sim\beta\to\delta\). Consequently an
\(\alpha\)-normal world satisfying \(\beta\) is \(\beta\)-normal.
If \(\alpha\leq\beta\leq\gamma\), then
\(\alpha\sim\gamma\to\beta\); hence an \(\alpha\)-normal world
satisfying \(\gamma\) also satisfies \(\beta\) and is \(\beta\)-normal.
These are the dependencies behind source Lemmas~22 and~25--27,
not presumed coherence properties of a newly named preorder.
\end{hblemma}

\begin{hbdefinition}{def:p08-klm-canonical}{The canonical copied-state structure}
For a preferential relation put
\[
 S=\{(m,\alpha):\alpha\in L,\ m\in N(\alpha)\},\quad
 l(m,\alpha)=m,
\]
and define
\[
 (m,\alpha)\prec(n,\beta)
 \quad\Longleftrightarrow\quad
 \alpha\leq\beta\ \text{ and }\ m\not\models\beta.
\]
This is a set construction even for infinite \(L,U\). The valuation
test concerns the first state's label; replacing it by a comparison
of labels or discarding the formula-coordinate copies changes the model.
\end{hbdefinition}

\begin{hblemma}{thm:p08-klm-canonical-order}{proved=annex:p08.klm.5}{The constructed preference is a strict partial order}
The displayed \(\prec\) is irreflexive and transitive. Its transitivity
uses the normal-world transport lemma, not just transitivity of \(\leq\).
\end{hblemma}

\begin{hblemma}{thm:p08-klm-canonical-minima}{proved=annex:p08.klm.6}{Characterizing every canonical minimal state}
A canonical state \((m,\beta)\) is minimal in \(\widehat\alpha\)
iff \(m\models\alpha\) and \(\beta\leq\alpha\).
Both directions are proved, including the witness obtained when
\(\alpha\vee\beta\not\sim\beta\).
\end{hblemma}

\begin{hblemma}{thm:p08-klm-canonical-smooth}{proved=annex:p08.klm.7}{Smoothness on every definable set}
For every formula \(\alpha\), every canonical state in
\(\widehat\alpha\) is minimal there or has a minimal predecessor there.
If \((m,\beta)\) is not minimal, a suitable normal world for
\(\alpha\vee\beta\) gives the explicit predecessor
\((n,\alpha\vee\beta)\). No well-foundedness or finiteness of
the state set is asserted or required.
\end{hblemma}

\begin{hbtheorem}{thm:p08-klm-pref-representation}{proved=annex:p08.klm.8}{General preferential representation, in both directions}
For the fixed compact \((L,U)\), a relation satisfies System P iff
it is exactly the relation of a smooth preferential model over \(U\).
The canonical copied-state structure above is such a model and satisfies
\(\alpha\sim\beta\) iff \(\alpha\sim_W\beta\) for all formulas.
This represents arbitrary set-sized relations and languages, not just
finite knowledge bases or finite sets of valuations. It is source
Theorem~3, with all its construction lemmas supplied here.
\end{hbtheorem}

\begin{hbcorollary}{thm:p08-klm-kb-completeness}{proved=annex:p08.klm.9}{Conditional knowledge-base completeness}
For any \(K\subseteq L\times L\), finite or infinite,
\(K\vdash_P\alpha\sim\beta\) iff every smooth preferential model
over \(U\) satisfying all assertions in \(K\) satisfies
\(\alpha\sim\beta\). Nonderivability supplies a canonical model
of the least P-closed relation containing \(K\), with an actual minimal
state falsifying the requested consequent. This is not a finite-model
claim for arbitrary infinite \(K\).
\end{hbcorollary}

\begin{hbtheorem}{thm:p08-klm-penguins}{proved=annex:p08.klm.10}{The penguin triangle and a finite exception}
From the source's defeasible assertions
\(p\sim b,\ p\sim\neg f,\ b\sim f\), P derives
\[
 p\wedge b\sim\neg f,\quad f\sim\neg p,\quad b\sim\neg p,
 \quad b\vee p\sim f,\quad b\vee p\sim\neg p.
\]
A two-state model over all valuations on \(b,p,f\) has
\(s_0\prec s_1\), with labels
\((b,p,f)(s_0)=(1,0,1)\),
\((b,p,f)(s_1)=(1,1,0)\).
It satisfies those original defaults, refutes \(p\sim f\), and
shows failure of arbitrary left strengthening from \(b\sim f\)
to \(b\wedge p\sim f\). Yet adding the expected \(f\) to \(b\)
preserves its minimal states and all its consequences.
The default \(p\sim b\) has not been replaced by a strict
background axiom \(p\to b\).
\end{hbtheorem}

\begin{hbtheorem}{thm:p08-klm-diamond}{proved=annex:p08.klm.11}{The source diamond: entailments and countermodels}
Use the full valuation universe on the listed atoms, and let
\(K=\{t\sim p,t\sim s,p\sim e,s\sim\neg e\}\).
Then P derives \(\top\sim\neg t\) and
\(\top\sim\neg(p\wedge s)\), but neither \(t\sim e\) nor
\(t\sim\neg e\) follows. Nor do \(p\sim\neg s\),
\(s\sim\neg p\), or \(a\wedge p\sim e\) for a fresh atom \(a\).
The annex gives concrete finite smooth models for each failure, not
a selection between presumed default extensions.
\end{hbtheorem}

\begin{hbtheorem}{thm:p08-klm-copies}{proved=annex:p08.klm.12}{Copies can be indispensable}
Over the two-atom language \(p,q\), take four states with
\(s_0\prec s_2\), \(s_1\prec s_3\), and no other comparisons.
Their labels are respectively \(p\neg q,\neg p\neg q,pq,pq\).
This finite smooth model has no equivalent smooth model over the same
four-valuation universe whose labeling function is injective. The two
copies of \(pq\) cannot simply be merged. This supplies the example
whose separating proof is left as an exercise in source \S5.2, p.~28.
\end{hbtheorem}

\begin{hbtheorem}{thm:p08-klm-nonsmooth}{proved=annex:p08.klm.13}{Minimal existence is weaker than smoothness}
There is an infinite strict-partial-order model over two atoms \(p,q\)
in which every nonempty formula-definable set has a minimal state,
but Cautious Monotonicity fails.
Take states \(a,d,b_0,b_1,\ldots\), comparisons
\(a\prec d\) and \(b_j\prec b_i\) iff \(j>i\), and no others;
label \(a\) by \((p,q)=(1,1)\), \(b_0\) by \((1,0)\), and
\(d,b_i\ (i\geq1)\) by \((0,0)\).
Then \(\top\sim_W p\) and \(\top\sim_W q\), but
\(p\not\sim_W q\). This is a concrete failure without smoothness,
even though merely asking for a minimal state in each definable set succeeds.
\end{hbtheorem}

\begin{hbtheorem}{thm:p08-klm-not-rational}{proved=annex:p08.klm.14}{Preferential need not be ranked or rational}
Take three states with labels \((p,q)=(0,1),(1,1),(1,0)\), and
only the first state preferred to the third. It is finite and smooth,
so validates P, but has \(\top\sim_W q\),
\(\top\not\sim_W\neg p\), and \(p\not\sim_W q\).
Thus Rational Monotony is not a P rule. This order is not ranked.
The ranked representation and rational-closure treatment belong to the
separate Lehmann--Magidor presentation, not to the P construction here
\cite{LehmannMagidor1992}.
\end{hbtheorem}

\begin{hbtheorem}{thm:p08-klm-cumulative-separation}{proved=annex:p08.klm.15}{Why cumulative states and preferential states differ}
The source's cumulative models label a state by a nonempty set of worlds
and make it satisfy \(\alpha\) when every label-world satisfies it.
In the one-state, empty-preference model labeled by the two valuations
of a single atom \(p\), all five C rules hold, but Or fails:
\(p\sim\bot\) and \(\neg p\sim\bot\), while
\(p\vee\neg p\not\sim\bot\). This distinguishes the cumulative
system from P; the singleton-world labeling condition is substantive.
\end{hbtheorem}

The representation above is for flat conditional assertions over the
specified classical universe. It does not import nested counterfactual
connectives, probability thresholds, AGM revision postulates, or the
separate representation constructions for C, CL, CM, M and ranked rational
relations. Those systems have different hypotheses and source scopes.
